\documentclass[12pt, a4paper]{article}
\usepackage{graphicx} % Required for inserting images
\usepackage{amsmath}
\usepackage{amsthm}
\usepackage{amssymb}
\usepackage[utf8]{inputenc}
\usepackage[left=2cm, top=-5cm, right=2cm, bottom=2cm, includeheadfoot, headheight=130pt, a4paper]{geometry}
\usepackage[polish]{babel}
\usepackage[utf8]{inputenc}
\usepackage{polski}
\usepackage[T1]{fontenc}
\usepackage{longtable}
\usepackage{caption}
\usepackage{subcaption}
\usepackage{booktabs}
\usepackage{array}
\usepackage{etoolbox}

\BeforeBeginEnvironment{verbatim}{\small}

\title{Wybrane Zagadnienia Algebry - Sprawozdanie 1}
\author{Rafał Włodarczyk 279762}
\date{Kwiecień 2026}

\begin{document}

\maketitle

\section{Zadanie 1 - Pierścień Liczb Gaussa}

Pierścień Liczb Gaussa ($\mathbb{Z}[i]$ - \textit{Gaussian Integers}) to zbiór liczb:
\begin{align*}
    \mathbb{Z}[i] = \{ a + bi : a, b \in \mathbb{Z} \}, \quad i^2 = -1
\end{align*}
Wyróżniamy dwie składowe $z\in \mathbb{Z}[i]$ - rzeczywista $a$ i urojona $b$. Dla zadanych danych wystarczające jest zdefiniowanie
struktury danych, w oparciu o typ całkowity \texttt{int64\_t}:
\begin{verbatim}
class Gauss {
  int64_t real;
  int64_t imaginary;
}
\end{verbatim}
Klasa \texttt{Gauss} zawiera następujące metody:
\begin{itemize}
    \item Norma zgodna z definicją $N(a+bi) = a^2 + b^2$. Dla zadanych danych $a+bi = 2+7i$:
    \begin{verbatim}
Norm(2+7i)
-> N = 53
    \end{verbatim}
    \item Dzielenie z resztą $X/Y$: $X=QY+R$, gdzie $N(R) < N(Y)$. Dzielenie odbywa się w $\mathbb{C}$, a następnie jest zaokrąglane do najbliższych liczb całkowitych.
    \begin{align*}
        Q &= \frac{a+bi}{c+di} \cdot \frac{c-di}{c-di} = \left\lfloor \frac{ac+bd}{N(y)} \right\rceil + i \left\lfloor \frac{bc-ad}{N(y)} \right\rceil\\
        R &= X - QY
    \end{align*}
    Na podstawie funkcji $\lfloor \cdot \rceil$ (zaokrąglenie do najbliższej liczby całkowitej)
    można zdefiniować 4 przypadki dzielenia, poprawne jeśli $N(R) < N(Y)$.
    Dla zadanych danych $X = (c+a) + (d+b)i = 11 + 14i$ oraz $Y = 6 + 2i$:
    \begin{verbatim}
Div(11+14i,6+2i), N(Y)=40
-> Q = 2+1i ; R = 1+4i, N(R)=17
-> Q = 2+2i ; R = 3-2i, N(R)=13
-> Q = 3+1i ; R = -5+2i, N(R)=29
-> Q = 3+2i ; R = -3-4i, N(R)=25
    \end{verbatim}
    Co zgadza się z oczekiwaniami ponieważ $N(R) < N(Y)$ dla wszystkich 4 przypadków oraz w $\mathbb{C}$ wynik dzielenia jest równy $2.35 + 1.55i$,
    widzimy jednocześnie że naturalne zaokrąglenie (od $0.5$ w górę) daje najmniejszą normę reszty - wynik najbliższy temu w liczbach zespolonych.
    \item Wyznaczenie NWD oraz NWW. NWD jest wyznaczany za pomocą standardowego Algorytmu Euklidesa, wykorzystującego wyżej zdefiniowane dzielenie z resztą.
    NWW jest wyznaczany na podstawie NWD, zgodnie z zależnością $LCM(X,Y) = \frac{XY}{GCD(X,Y)}$. Dla zadanej listy $(a+bi,c+di,e+di)$:
    \begin{verbatim}
Array=(2+7i,9+7i,6+7i)
-> GCD = 1+0i, -1+0i, 0+1i, 0-1i
-> LCM = -245-725i, 245+725i, -245+725i, 245-725i
    \end{verbatim}
    Jeśli $GCD(X,Y) = 1$, to każdy $ \textit{unit} \in\{1,-1,i,-i\}$ dzieli $X$ i $Y$ wobec tego znów mamy 4 przypadki NWD i NWW.
    Wywołania dla listy pustej oraz jednoelementowej dają odpowiednio:
    \begin{verbatim}
GCD () = 0+0i
GCD (2+7i) = 2+7i
    \end{verbatim}
\end{itemize}

\section{Zadanie 2}

Pierścień wielomianów $\mathbb{R}[x]$ to zbiór wielomianów o współczynnikach rzeczywistych:
\begin{align*}
    \mathbb{R}[x] = \{ a_n x^n + a_{n-1} x^{n-1} + ... + a_1 x + a_0 : n \in \mathbb{N}, a_i \in \mathbb{R} \}
\end{align*}
Reprezentacja wielomianów jest realizowana za pomocą wektora współczynników:
\begin{verbatim}
class Poly {
  std::vector<double> c;
}
\end{verbatim}
Gdzie \texttt{c[i]} to współczynnik przy $x^i$. Najprostszą optymalizacją pamięciową byłoby wykorzystanie struktury sparse vector,
przy czym nie jest ona potrzebna do wykonania obliczeń dla zadanych niewielkich danych.
Klasa \texttt{Poly} zawiera następujące metody:
\begin{itemize}
    \item Norma $N(w)=N(a_n x^n + \dots + a_0) = \max_i {i : a_i > 0} = \text{st}(w)$ - stopień wielomianu. Dla zadanych danych $(cx^a+b)$:
    \begin{verbatim}
Norm(9x^2 + 7)
2
    \end{verbatim}
    \item Dzielenie z resztą - algorytm dzielenia w słupku, gdzie dzielnik jest odejmowany od dzielnej dopóki norma reszty jest większa lub równa normie dzielnika.
    Dla zadanych danych $X = cx^a + b = 9x^2 + 7$ oraz $Y = x + 1$:
    \begin{verbatim}
Div(9x^2 + 7;x + 1)
Q = 9x - 9
R = 16
    \end{verbatim}
    \item NWD,NWW - analogicznie do zadania 1, z wykorzystaniem algorytmu Euklidesa oraz zależności $LCM(X,Y) = \frac{XY}{GCD(X,Y)}$.
    \item EE - rozszerzony algorytm Euklidesa - dla wielomianów $V$,$W$ znajduje $v$,$W$ takie że $vV + wW = GCD(v,W)$. 
    Dla zadanych danych $v(x)=ax^3+bx^2+cx+d$ oraz $w(x)=dx^3+ex^2+fx$ zachodzi:
    \begin{verbatim}
Extended GCD(2x^3 + 7x^2 + 9x + 7, 7x^3 + 6x^2 + 2x)
GCD = 1
v = 7.29x^2 - 45.55x + 30.63
w = -40.05x^3 + 471.35x^2 - 1550.08x + 928.84
    \end{verbatim}
    Na tej podstawie jesteśmy w stanie znaleźć $g : 1 \notin GCD(v(x),w(x)+g)$ - czyli $v(x), w(x)+g$ muszą współdzielić pierwiastek.
    Wyznaczmy numerycznie pierwiastek $v(\alpha)=0$ (metodą bisekcji), następnie wyznaczmy $g$ takie że $w(\alpha)+g=0$:
    \begin{align*}
        g = -w(\alpha)
    \end{align*}
    Dla zadanych danych:
    \begin{verbatim}
Root of v: -2.1693702696445882
w at root of v: -47.56767332520533
LCM(2x^3 + 7x^2 + 9x + 7, 7x^3 + 6x^2 + 2x + 47.57)
LCM = 14x^6 + 61x^5 + 109x^4 + 212.14x^3 + 392.98x^2 + 442.11x + 332.98
    \end{verbatim}
\end{itemize}

\section{Zadanie 3}

Porządek produktowy $\leq$ na $\mathbb{N}^n$ jest definiowany jako:
\begin{align*}
    (a_1, a_2, \dots, a_n) \leq (b_1, b_2, \dots, b_n) \iff \left(\forall i = 1,2,\dots, n \right) a_i \leq b_i 
\end{align*}
Wykonajmy sprawdzenie dla zadanych par $(a,b),(c,d),(e,f)=(2,7),(9,7),(6,2)$ (każdy z każdym):
\begin{verbatim}
compare([2, 7], [2, 7]) is ==
compare([2, 7], [9, 7]) is <=
compare([2, 7], [6, 2]) is no comparison possible
compare([9, 7], [2, 7]) is >=
compare([9, 7], [9, 7]) is ==
compare([9, 7], [6, 2]) is >=
compare([6, 2], [2, 7]) is no comparison possible
compare([6, 2], [9, 7]) is <=
compare([6, 2], [6, 2]) is ==
\end{verbatim}
Wykonajmy sprawdzenie dla trójek $(a,c,e),(b,d,f)=(2,9,6),(7,7,2)$:
\begin{verbatim}
compare([2, 9, 6], [7, 7, 2]) is no comparison possible
\end{verbatim}
Następnie tworzymy algorytm znajdujący dla danego zbioru punktów $A \subset \mathbb{N}^n$ zbiór elementów $\leq$-minimalnych:
\begin{align*}
    \text{Min}(A) = \{ a \in A : \nexists b \in A : b < a \}
\end{align*}
Algorytm naiwny sprawdza dla każdego punktu $a \in A$ czy istnieje punkt $b \in A$ taki że $b < a$. Jeśli tak, to $a$ nie jest $\leq$-minimalny.
Jego złożoność można ocenić na $\mathcal{O}(n^2)$, gdyż dla każdego z $n$ punktów sprawdzamy pozostałe $n-1$ punkty. 
Algorytm ten jest jednak zupełnie nieefektywny dla zadanego zestawu danych, dlatego proponowany algorytm wykorzystuje sprytniejsze sprawdzanie.
\begin{verbatim}
for (p in A) 
    is_min = true
    for (q in MIN)
        if (p < q)
            MIN.remove(q)
            is_min = false
            break
        else if (q < p)
            is_min = false
            break
    if (is_min)
        MIN.add(p)
\end{verbatim}
Musimy sprawdzić każdy element $p$ z $A$. Następnie budujemy zbiór $\text{MIN}$, który będzie zawierał $\leq$-minimalne elementy. 
Jeśli znajdziemy taki element $p$, że istnieje $q \in \text{MIN}$ taki że $p < q$, to usuwamy $q$ ze zbioru $\text{MIN}$, ponieważ $q$ nie jest już $\leq$-minimalny.
Jeśli zachodzi $q < p$ to na pewno $p$ nie jest $\leq$-minimalny, więc możemy przejeść dalej.
Złożoność tego algorytmu zależy od rozmiaru zbioru $\text{MIN}=m$ i wynosi $\mathcal{O}(nm)$. W najgorszym przypadku, gdy wszystkie elementy są $\leq$-minimalne, złożoność będzie $\mathcal{O}(n^2)$,
ale w praktyce jest znacznie lepsza, ponieważ $|\text{MIN}| \ll |A|$ dla zadanych danych. Szukamy elementów $\leq$-minimalnych w zbiorze
\begin{itemize}
    \item $A = \{(x,y)\in \mathbb{N}^2 : (x-a)^2 + (y-b)^2 < 5\}$, gdzie $a=2$ i $b=7$.
\begin{verbatim}
Count Min elements of A: 2
Min elements of A:
[0, 6], [1, 5]
\end{verbatim}
Następnie szukamy elementów $\leq$-minimalnych w zbiorze
    \item $B = \{(x_1,x_2,x_3,x_4) \in \mathbb{N}^4 : (x_1-c)^2 + (x_2-d)^2 + (x_3-e)^2 + (x_4-f)^2 > 224\}$, gdzie $c=9$, $d=7$, $e=6$ i $f=2$.
\end{itemize} 

\begin{verbatim}
Count Min elements of B: 30
Min elements of B:
[0, 0, 0, 10], [0, 0, 13, 9], [0, 0, 14, 8], [0, 0, 15, 6], [0, 0, 16, 0]
[0, 15, 0, 9], [0, 15, 13, 8], [0, 15, 14, 6], [0, 15, 15, 0], [0, 16, 0, 8]
[0, 16, 13, 6], [0, 16, 14, 0], [0, 17, 0, 5], [0, 17, 13, 0], [0, 18, 0, 0]
[19, 0, 0, 9], [19, 0, 13, 8], [19, 0, 14, 6], [19, 0, 15, 0], [19, 15, 0, 7]
[19, 15, 13, 6], [19, 15, 14, 0], [19, 16, 0, 5], [19, 16, 13, 0], [19, 17, 0, 0]
[20, 0, 0, 7], [20, 0, 13, 5], [20, 0, 14, 0], [20, 15, 0, 0], [21, 0, 0, 0]   
\end{verbatim}

\end{document}